\section{Limits and further questions}
The exact formulas \eqref{pdt:eq:T}--\eqref{pdt:eq:transform}, the ascent-sequence/table bridge and the historical rigid subclass belong to the cited enumerative literature. The contribution developed here is the relative saddle analysis with its explicit error organization and consequences for this specific model. The result is self-contained modulo standard gamma asymptotics and Gaussian Poisson summation, whose needed uniform estimates have been explained.

Several useful questions remain separate from the results proved here.
\begin{enumerate}
\item \textbf{Effective error constants.} Explicit uniform gamma bounds and validated saddle evaluation could turn the eventual inverse brackets into finite-size certificates. Our finite diagnostics do not accomplish this.
\item \textbf{Growing expansion order.} The coefficient algorithms determine every fixed order. Their growth as the order increases, optimal truncation and any exponentially improved remainder require additional analysis.
\item \textbf{Complex dimension marks.} Bounded real marks suffice for the dimension laws and moments established here. Uniform complex marks or a global analytic generating-function theorem are not asserted.
\item \textbf{Further statistics.} The expectations in \eqref{pdt:eq:bad-diag}--\eqref{pdt:eq:bad-coll} suggest studying the joint collision structure. A full distribution theorem, particularly one strong enough for relative rare probabilities, is not a consequence of those expectations.
\item \textbf{Other difference parameters.} The $d=0$ and $d=1$ exact models admit the present analysis. General fixed positive $d$, and the stabilized $d\to\infty$ family, must be treated from their own enumeration and cannot be inferred by changing the signs in our formulas.
\end{enumerate}
The asymptotic conclusions do not establish worldwide priority, absence of equivalent results in uninspected sources, external peer review or machine-checked proofs. The accompanying source and computation notes delimit what was actually checked. In particular, source text access is not represented as inspection of unavailable page images, and finite numerical agreement is not represented as a theorem about an infinite range.

\subsection{Further questions and research}\label{pdt:sec:further}
\begin{writenote}[Added 5 October 2026, batch~102]
Following Vladimir Reshetnikov's standing rule of 4~October 2026, every
claim of the source that is not proved in full is stated here as an open
question, with its source, the argument offered and what is missing.
Questions~\ref{pdt:q:effective}--\ref{pdt:q:difference} are the five items
above; Questions~\ref{pdt:q:moments}--\ref{pdt:q:moderate} are steps that
the text asserts or condenses; Question~\ref{pdt:q:literature} concerns the
external inputs. The intake found no false claim in the source; the one
point of wording, in \eqref{pdt:eq:inverse-separation}, is explained in the
note after Theorem~\ref{pdt:thm:seeds}.
\end{writenote}
\begin{enumerate}
\item\label{pdt:q:effective} \textbf{Effective constants and certified
 thresholds} (source, item~1 above). Every $O$-constant and threshold in
 Theorems~\ref{pdt:thm:main}--\ref{pdt:thm:seeds} is an existence constant
 (Remark~4.2). Explicit uniform gamma and polygamma bounds in
 Lemma~\ref{pdt:lem:global}, \eqref{pdt:eq:derivativecompare} and
 Lemma~\ref{pdt:lem:poisson}, with validated evaluation of the saddle,
 would turn the brackets \eqref{pdt:eq:inversebracket},
 \eqref{pdt:eq:bseedbracket} and \eqref{pdt:eq:cseedbracket} into finite
 certificates. The finite evidence shows why this matters: $C_1$ changes
 sign only near $N\approx1.59\times10^4$ (Remark~\ref{pdt:rem:C1sign}), the
 displayed limits are approached slowly (Remark~\ref{pdt:rem:slowbinary}),
 and the elementary seed rounds the wrong way at $X=b_{1000}$
 (Section~\ref{pdt:sec:inverse}). An interval computation of
 \eqref{pdt:eq:C1} would also make the sign change of
 Remark~\ref{pdt:rem:C1sign} a theorem.
\item\label{pdt:q:order} \textbf{Growing expansion order} (source,
 item~2). How do $C_\ell$ and $Q_k$ grow with the order; is there an
 optimal truncation, and an exponentially improved remainder? Missing:
 any estimate uniform in the order. The source's data at $N=1000$
 (Section~\ref{pdt:sec:checks}): the residuals after $K=1,2,3$ terms are
 about $-1.13\times10^{-5}$, $6.07\times10^{-10}$ and
 $-5.56\times10^{-13}$, gains of factors about $1.9\times10^4$ and
 $1.1\times10^3$.
\item\label{pdt:q:complex} \textbf{Complex dimension marks} (source,
 item~3). Theorem~\ref{pdt:thm:main} is uniform for bounded real $t$
 only. Uniform complex marks, or a global analytic statement about
 $B(z,h)$, would give, for example, large-deviation rates and analytic
 continuation in the mark; nothing of this is asserted.
\item\label{pdt:q:collision} \textbf{The joint collision structure}
 (source, item~4, and Section~8.2). The expectations
 \eqref{pdt:eq:bad-diag}--\eqref{pdt:eq:bad-coll}, asymptotically $u$ and
 $u^2/2$, ``explain'' the exponent of $p_N\sim e^{-u^2/2-u}$, as if
 $(D_{\rm bad},C_{\rm bad})$ were nearly independent Poisson variables.
 The source states this as an interpretation only. Missing: a joint
 distribution theorem for the two bad counts that is strong enough to give
 the relative probability of the rare event $\{D_{\rm bad}=C_{\rm bad}=0\}$
 (a total-variation statement is not).
\item\label{pdt:q:difference} \textbf{Other difference parameters}
 (source, item~5). Only $d=0$ and $d=1$ of~\cite{ccs} are treated. Fixed
 $d\ge2$ and the stabilized $d\to\infty$ family need their own exact
 enumeration and analysis; changing signs in the present formulas proves
 nothing.
\item\label{pdt:q:moments} \textbf{Higher moments and local corrections}
 (source, the paragraph after Proposition~\ref{pdt:prop:moments}). The
 source asserts that ``all higher fixed moments and local corrections are
 obtainable from the weighted-contraction rule and finite series
 division''. Sketch: repeat Section~\ref{pdt:sec:transfer} with a factor
 $z^k$, which the polynomial tail bound after \eqref{pdt:eq:relative-tail}
 makes uniformly integrable. Missing: the general statement with its
 remainder (only $k=1,2$ are proved, in
 Proposition~\ref{pdt:prop:moments}) and the local expansion
 \eqref{pdt:eq:localexpansion} beyond its first correction.
\item\label{pdt:q:smooth} \textbf{The smooth elementary carrier}
 (source, proof of Theorem~\ref{pdt:thm:seeds}). The real-size estimate
 \eqref{pdt:eq:smooth-elementary},
 $F_K(y)=F_0(y)+g_+(u(y))+O_K(u^4/y)$, is asserted as ``the smooth
 derivation'' of Theorem~\ref{pdt:thm:elementary}, and the bound
 $C_\ell'(y)=O_\ell(y^{-\ell-1})$ in \eqref{pdt:eq:coeffderivatives} is
 argued in one sentence. Sketch: the proof of
 Theorem~\ref{pdt:thm:elementary} uses only the exact phase at real
 arguments, and each contraction is a polynomial in the $\lambda_j$, whose
 $y$-derivatives \eqref{pdt:eq:coeffderivatives} follow from
 \eqref{pdt:eq:mixedbounds} and \eqref{pdt:eq:rderivative}. Missing: the
 written argument that every remainder in the proof of
 Theorem~\ref{pdt:thm:elementary} is uniform for real $y$, and the
 chain-rule bookkeeping for the contractions.
\item\label{pdt:q:binaryremainder} \textbf{The remainder in
 Theorem~\ref{pdt:thm:binary-all}} (source, its proof). The statement
 that ``higher remainders are bounded by the same order after integration
 against a slightly weakened Gaussian'' is a sketch. Missing: the explicit
 domination of the Taylor remainders of $f_N$ and $h_N$ on the window
 $|z|\le N^\varepsilon$ that yields $L^{4K+4}N^{-K-1}$ without additional
 logarithmic factors. The intake checked the window and tail argument and
 believes the bound.
\item\label{pdt:q:moderate} \textbf{Moderate deviations} (source, after
 Proposition~\ref{pdt:prop:moments}): the relative Gaussian approximation
 for $|z|=o(N^{1/6})$, argued in two sentences. \emph{Resolved at the
 write:} Proposition~\ref{pdt:prop:moderate} proves it, with the error
 $\exp(O(|z|^3N^{-1/2}+N^{-1}))$. Listed for the record.
\item\label{pdt:q:literature} \textbf{External inputs and priority}
 (source, Section~\ref{pdt:sec:models} and \nolinkurl{SOURCES.md}). The
 source's search was bounded, and its page-level readings of
 \cite{bek,khamis,bcdk,dm,ccs,kr,bcuf,cfz,hj} were not checked by ProveIt.
 Open: whether the relative results here, in particular the answer to
 A098568's question, appear in sources the search did not reach, and
 whether Kote\v{s}ovec's finer logarithmic expressions~\cite{oeis} have a
 published proof (the inspected records cite none).
\end{enumerate}

\paragraph{Independent check of the write (dated note, 5~October 2026).}
An independent adversarial check of the three results this write supplied,
made by the intake after the write (\nolinkurl{3daaab24e}), found all three
valid: Proposition~\ref{pdt:prop:moderate}, the window of the note after
Theorem~\ref{pdt:thm:seeds}, and Remark~\ref{pdt:rem:C1sign}. It found no
counterexample and no gap in any proof chain. Two changes followed. The
note before Proposition~\ref{pdt:prop:moderate} listed the inputs of the
proof without \eqref{pdt:eq:rscale}, which the proof uses in its first
sentence; the list now includes it. In Remark~\ref{pdt:rem:C1sign} the
statement that $C_1<0$ until $N$ lies ``between $10^4$ and $10^6$'' was true
but loose; it now reads ``between $10^4$ and $10^5$'', with a dated note
keeping the first wording. The check reread every step and used neither the
delivered programs nor the write's. Its own programs (Python~3.14.4,
mpmath~1.3.0) evaluate the exact gamma phase \eqref{pdt:eq:phase}, with
derivatives by the Fa\`a di Bruno formula for a composition with a
quadratic (checked against numerical differentiation to $10^{-49}$), find
the saddle by sign-checked bisection on $I_N$, and compute $C_\ell$ by
direct enumeration of the contraction \eqref{pdt:eq:contraction}, which
agrees exactly with \eqref{pdt:eq:C1}.

For Proposition~\ref{pdt:prop:moderate}: the range $|m-r|\le r/2$ lies in
$I_N$ for large $N$ (for $T=1$ from about $L>14$ on), the cubic bound, the
Lagrange remainder, the normalization by \eqref{pdt:eq:exactshape} with
$K=1$ and the second assertion are correct, and the window
$|z|=o(N^{1/6})$ is the natural one, since $\lambda_3\sim\sqrt2\,N^{-1/2}$.
As a stress test, for $t=-1,0,1$ and $N=10^3,10^4,10^5$ the check computed
the exact probabilities $\Prob_t(D_N=m)$ (normalized over $r/4\le m\le3r$;
the omitted terms have relative weight below $e^{-1300}$) and took the
supremum, over the integers $m$ with $|m-r|\le r/2$ (every seventh one at
$N=10^5$), that is up to $|z|\approx250$, of
$\bigl|\log\bigl(\Prob_t(D_N=m)/(\sqrt A\,\varphi(z))\bigr)\bigr|$ divided by
$|z|^3N^{-1/2}+N^{-1}$. It is $0.24$--$0.26$ at $N=10^3$, $0.28$--$0.29$ at
$N=10^4$ and $0.30$--$0.31$ at $N=10^5$, attained at the left end
$m\approx r/2$. The slow growth matches the leading-order value
$4\sqrt2\,(\log2-\frac58)\approx0.385$ of this ratio at $m=r/2$, from
$\rho\approx2N(\log(1+u)-u+u^2/2)$ with $u=(m-r)/r$; no failure of
uniformity appeared.

For the note after Theorem~\ref{pdt:thm:seeds}: the window was rederived
from the brackets \eqref{pdt:eq:bseedbracket}--\eqref{pdt:eq:cseedbracket}
and the exact gap $z_c-z_b=v/2+1$. With exact $b_n$ and $c_n$ for
$n\le600$ (big-integer sums), for 4000 log-uniformly sampled $X$ in
$[b_{60},c_{600}]$ both brackets held with $C=1$, and
$\nu_c(X)-\nu_b(X)-v/2$ ranged over $[0.035,1.764]$, inside the window. The
seed errors $z_b(b_n)-n$ and $z_c(c_n)-n$ are at most $0.06\,v^3/n_0$ in
absolute value for $n=50,\ldots,3200$, as \eqref{pdt:eq:seederror} allows.
A side finding: $b_n$ is strictly increasing on $1\le n\le600$, but $c_n$
increases only weakly ($c_0=c_1=c_2=1$), consistent with the source, whose
injection gives $c_{n+1}\ge c_n$ and which needs no more.

For Remark~\ref{pdt:rem:C1sign}: the check reproduced every entry of the
table; on the 51-point grid the last negative point is $N=15{,}849$ and
the first positive one $N=19{,}953$, with least value $-0.0217$ at $N=50$
(over the integers $10\le N\le199$ the least value is $-0.021713$, at
$N=48$); $NC_1=-1.0722\times10^{-7}$ at $N=15{,}864$ and
$+1.0975\times10^{-7}$ at $N=15{,}865$; and a 600-point logarithmic grid on
$[10,10^6]$ shows exactly one sign change. The heuristic
$-\frac38+\frac5{12}=\frac1{24}$ and $1/L\approx0.07$ at $N=10^6$ are
correct. For $N\le46$ the interval $I_N$ does not bracket the saddle
(there $4N/L>N$), so the check took the root on $[1,N)$; this affects only
grid points below $50$.

Spot checks of the source: $b_1,\ldots,b_9=1,2,5,14,43,143,510,1936,7775$,
the row $(1,6,6,1)$ at $N=4$ and the binomial transform
\eqref{pdt:eq:bin-transform} for $N\le30$; at $N=1000$ the saddle, mean,
variance, $p_N$ and the residuals for $K=1,2,3$ of
Section~\ref{pdt:sec:checks}, and the write's corrected-moment errors
$7.0\times10^{-6}$ and $1.47\times10^{-6}$, all to every printed digit; and
the rounding counterexample at $X=b_{1000}$: $z_b=1000.0014164605$, whose
ceiling $1001$ exceeds the threshold $1000$ ($b_{999}<b_{1000}$), the $K=2$
smooth root $1000.00000000014117$ by the check's own bisection, and
$z_c(c_{1000})=1000.0028242626$.

Separately, a parallel check of Report~243's write found that this
report's credit to that report (Section~\ref{pdt:sec:provenance}) said
that the estimate of its Lemma~2.1 follows from
Theorem~\ref{pdt:thm:elementary}; that holds for the estimate of $b_n$
but not for the lemma's single-term statement, which is the display after
\eqref{pdt:eq:fixedalpha}. The credit is corrected where it stands, with a
dated note keeping the first wording.

This was a careful reading with numerical tests in 30- to 60-digit floating
point, not interval-certified, and not a formal verification or an
external review.
